\chapter{言語モデルの学習と計算}\label{ch:llm}

前の章では、できあがった言語モデルを「使う」方法を説明した。この章では、言語モデルに\textbf{学習させる}方法（追加学習、LoRA）と、それを動かす計算（GPU、数の精度）を説明する。実験 4 と 10 で何をしたかを理解するための章である。

\section{学習の 3 つの段階}

\begin{longtable}{P{3.2cm}P{5.6cm}P{5.6cm}}
\toprule
\textbf{段階} & \textbf{中身} & \textbf{この研究} \\
\midrule
\endhead
\bottomrule
\endfoot
事前学習 & 大量の文章で、次のトークン（または隠した語）を当てる。数十億〜数兆のパラメータを、膨大な計算で学習する & 行っていない。公開されたモデル（gpt-4o-mini、ChronoGPT、FinBERT）を使った \\
追加学習（ファインチューニング） & 事前学習したモデルを、個別の課題の正解つきデータで少しだけ学習し直す & FinBERT の全体（実験 4・11）、ChronoGPT の LoRA（実験 10） \\
文脈の中の学習 & 重みは変えず、プロンプトに例や情報を入れて振る舞いを変える & H07（自分の過去の成績をプロンプトで渡す）。これは「学習」ではない \\
\end{longtable}

\section{全体の追加学習}

モデルのすべての重みを、課題のデータで学習し直す。FinBERT（約 1.1 億パラメータ）なら、この PC の GPU で十分にできる。実験 4 では、分類の頭（3 つのラベル）を回帰の頭（数 1 つ）に替え、決算発表文の先頭 512 トークンから 2 取引日の反応を当てるように、2015〜2022 年の 13,979 本で 2 エポック学習した（1,748 歩、756 秒）。

大きなモデル（数十億パラメータ）では、全体の追加学習は難しい。重みと、その勾配と、最適化の状態（Adam は重み 1 つにつき 2 つの数を持つ）を GPU のメモリに置くと、重みの数の十数倍のバイトが要る。ChronoGPT（1.55B）を全体で学習するには数十 GB のメモリが要り、この PC の 8GB では足りない。

\section{パラメータ効率のよい追加学習（PEFT）}

元の重みは固定し、少数の新しい重みだけを学習する方法を、まとめて \term{PEFT} と呼ぶ。

\subsection{LoRA}

\term{LoRA}（低ランク適応）は、重みの行列 $W$（$d\times d$）の更新を、2 つの細長い行列の積で表す。
\begin{equation}
W' = W + \frac{\alpha}{r}\,BA,\qquad B\in\mathbb{R}^{d\times r},\; A\in\mathbb{R}^{r\times d},\; r\ll d
\end{equation}
$W$ は固定し、$A$ と $B$ だけを学習する。$r=8$、$d=1536$ なら、1 つの行列あたり $2\times8\times1536=24{,}576$ 個で、元の $1536^2\approx236$ 万個の約 1\% である。$B$ を 0 から始めるので、学習の前のモデルは元のモデルと完全に同じである。

\begin{iobox}[入力と出力：実験 10 の LoRA]
\iorow{元のモデル}{ChronoGPT-Instruct の版 2022。重みはすべて固定}
\iorow{足した重み}{52 層すべての注意の \code{c\_q} と \code{c\_v} に、ランク 8、$\alpha$ 16 の LoRA。最後の層の平均（1,536 個）に線形の頭。合計 約 256 万（モデル全体の約 0.17\%）}
\iorow{入力}{決算発表文の先頭 1,792 トークン}
\iorow{正解}{2 取引日の反応（学習の標本の 1\%・99\% 点で丸めて標準化）}
\iorow{学習}{2015〜2022 年の 13,979 本、1 エポック、実質のバッチ 16（873 歩）。AdamW（LoRA 2e-4、頭 1e-3）、最初の 10\% で学習率を上げて直線で下げる。勾配チェックポイントでメモリを節約。7.6 時間、GPU のメモリの最大 5.63GB}
\iorow{出力}{2023 年以降の 7,524 本の反応の予測（75 分）}
\iorow{結果}{損失は約 1.0（何も学ばないときの値）からほとんど下がらなかった。予測は +29.5bp（$t$ 2.40）だが、追加学習した FinBERT と同時で +0.4bp、実験 5 のリッジと同時で $-1.3$bp}
\end{iobox}

\subsection{ソフトプロンプト}

\term{ソフトプロンプト}（プロンプト・チューニング）は、入力の先頭に「学習できる埋め込みのベクトル」を数個〜数十個付け、それだけを学習する。人が書くプロンプトの代わりに、勾配降下法で最適な「プロンプト」を探すと考えればよい。学習する数は LoRA よりさらに少ない。この研究では試していない。

\subsection{これは深層学習か}

LoRA もソフトプロンプトも、損失の勾配をネットワーク全体に逆向きに流し（誤差逆伝播）、重みを勾配降下法で更新する。動かす重みが少ないだけで、やっていることは深層学習そのものである。一方、実験 5 の「埋め込み＋リッジ」は、LLM の中は動かさず、出てきたベクトルに線形の回帰を置くだけなので、深層学習の部分は事前学習だけである（\term{線形プローブ}と呼ぶ）。

\section{なぜ LoRA はリッジに負けたか}

実験 10 の予想は「中まで学習させれば、リッジより強くなる」だったが、外れた。考えられる理由は次のとおりである。

\begin{enumerate}
\item \textbf{目的変数の雑音が大きい}：決算の反応のうち文章で説明できるのは $R^2$ で数 \% である。13,979 本で、約 256 万の重みを、この弱い信号に合わせるのは難しい。
\item \textbf{1 エポックで、設定を選び直していない}：学習率やエポック数を選び直すと過剰適合になるので、1 つの設定に固定した。もっと良い設定はありうる。
\item \textbf{リッジは強い基準}：線形プローブは、事前学習の知識を壊さずに使う。少ないデータでは、全体を動かすより線形の頭だけを学習するほうが強いことがよくある。
\end{enumerate}

FinBERT の全体の追加学習（実験 4）が効いたのは、モデルが小さく（1.1 億）、学習する数とデータの数の釣り合いが良かったためと考えられる。

\section{数の精度とメモリ}

\begin{longtable}{P{2.4cm}P{2.0cm}P{10.0cm}}
\toprule
\textbf{形式} & \textbf{1 つの大きさ} & \textbf{性質} \\
\midrule
\endhead
\bottomrule
\endfoot
fp32 & 4 バイト & 普通の浮動小数点。ChronoGPT の公開された重み（7.4GB） \\
bf16 & 2 バイト & 指数部は fp32 と同じ 8 ビットで、仮数部を 7 ビットに減らしたもの。値の範囲は fp32 と同じで、精度だけが下がる。深層学習の標準。この研究では ChronoGPT の重みを bf16 にして保存した（3.7GB） \\
fp16 & 2 バイト & 指数部が 5 ビットで、値の範囲が狭く、あふれやすい \\
int8・int4 & 1・0.5 バイト & \term{量子化}。重みを整数に丸め、メモリを大きく減らす。丸めの誤差で、出力が少し変わる \\
\end{longtable}

\begin{notebox}[bf16 は量子化ではない]
bf16 にしても、重みは浮動小数点のままで、値の範囲も同じである。bf16 は深層学習の学習と推論で標準的に使われる精度で、埋め込みのような出力への影響は小さい（この研究では、元の fp32 との差は測っていない）。int8・int4 の量子化は、出力に小さな雑音を足すことになり、反応という弱い信号を測る研究では避けたほうがよい。この研究では量子化は使っていない。
\end{notebox}

\section{計算の場所}

\begin{description}[leftmargin=1.5em,style=nextline,itemsep=0.4em]
\item[この PC の GPU] NVIDIA GeForce RTX 5050 Laptop（8GB）。CUDA 12.8 以上が要る。ChronoGPT の埋め込みは 1 本あたり約 0.4 秒で、実験 5 の約 3.7 万本に約 4.3 時間かかった。LoRA の学習は 1 本あたり 1.6〜1.9 秒。実験のほぼすべてをここで計算した。
\item[借りる GPU（RunPod）] 時間単位で GPU を借りるサービス。この PC で約 14 時間かかる見込みだった実験 11 の点数付けを短くするために使ったが、結果をダウンロードする前に Pod を止めてしまい、ディスクごと消えた。結局この PC で約 19 時間かけてやり直した。料金は Pod が動いている時間（設定や待ちの時間を含む）にかかる。借りた GPU は、この PC の 3 分の 1 程度の時間で済む見込みだったが、1 時間あたりの値段は 10 倍以上だった。メモリが足りない場合のほかは、得にならないことが多い。
\item[内蔵 GPU] CPU に内蔵された GPU は、CUDA が使えず、メモリも共有なので、PyTorch のこれらのモデルの計算にはほぼ使えない。
\end{description}

\section{RAG}

\term{RAG}（検索で補う生成）は、質問に関係する文書を検索して、プロンプトに入れてから LLM に答えさせる方法である。LLM の知識の締め切りより新しい情報や、社内の文書を使わせるのに向く。重みは変えない。この研究では、決算発表文そのものを直接プロンプトに入れたので、検索の部分は要らなかった。

\begin{projbox}
\begin{itemize}
\item \textbf{全体の追加学習}：FinBERT（実験 4、+71bp 単独）。実験 11 でも、切り方ごとに学び直した BERT が NLP と ChronoGPT より強かった（$R^2$ 増分 +1.01）。
\item \textbf{LoRA}：ChronoGPT 版 2022（実験 10）。リッジより弱い。
\item \textbf{線形プローブ}：ChronoGPT の埋め込み＋リッジ（実験 5・9・11）。LLM の中は動かさない。
\item \textbf{bf16}：ChronoGPT の重みの保存と計算。量子化はしていない。
\item \textbf{計算}：この PC の RTX 5050（8GB）。RunPod の結果は失われた。
\end{itemize}
\end{projbox}
